\documentclass[11pt,a4paper]{article}
\usepackage[utf8]{inputenc}
\usepackage[T1]{fontenc}
\usepackage[spanish]{babel}
\usepackage{amsmath,amssymb,booktabs}
\usepackage[margin=2.5cm]{geometry}

\title{Modelo 3-DOF consistente con \texttt{vrx::SimpleHydrodynamics}}
\author{}
\date{}

\begin{document}
\maketitle

\section{Convenciones}
El modelo se expresa en el marco solidario con el vehículo \texttt{base\_link}:
\begin{equation}
\nu = [u\;v\;r]^\top,
\end{equation}
donde $u$ es surge (eje $+x$), $v$ es sway (eje $+y$) y $r$ es la velocidad de guiñada. Los cuatro propulsores apuntan en $+x$ y se sitúan en $y=\pm0.29\,\mathrm{m}$. Por tanto,
\begin{equation}
 T_u=\sum_{i=1}^{4}T_i,\qquad
 T_r=-\sum_{i=1}^{4}y_iT_i.
\end{equation}

\section{Modelo estándar de Fossen}
El modelo planar 3-DOF estándar se escribe como
\begin{equation}
 M\dot\nu+C(\nu)\nu+D(\nu)\nu=\tau,
\end{equation}
con $M=\mathrm{diag}(m_{11},m_{22},m_{33})$, $\tau=[T_u\;0\;T_r]^\top$ y
\begin{equation}
C(\nu)=
\begin{bmatrix}
0&0&-m_{22}v\\
0&0&m_{11}u\\
m_{22}v&-m_{11}u&0
\end{bmatrix}.
\end{equation}
Con amortiguamiento diagonal cuadrático, sus ecuaciones escalares son
\begin{align}
m_{11}\dot u-m_{22}vr+X_u u+X_{uu}|u|u &=T_u,\\
m_{22}\dot v+m_{11}ur+Y_v v+Y_{vv}|v|v &=0,\\
m_{33}\dot r+(m_{22}-m_{11})uv+N_r r+N_{rr}|r|r&=T_r.
\end{align}

\section{Modelo implementado por el simulador tras la corrección}
El plugin ahora aplica la fuerza de Coriolis de masa añadida estándar
\begin{equation}
-C_A(\nu)\nu=
\begin{bmatrix}
 Y_{\dot v}vr\\
 -X_{\dot u}ur\\
 (X_{\dot u}-Y_{\dot v})uv
\end{bmatrix},
\end{equation}
además de $-M_A\dot\nu-D(\nu)\nu$. Gazebo ya integra la dinámica rígida; por ello, al sumar ambas contribuciones se obtiene exactamente:
\begin{align}
 (m+X_{\dot u})\dot u -(m+Y_{\dot v})vr
   +X_u u+X_{uu}|u|u &= T_u, \\
 (m+Y_{\dot v})\dot v +(m+X_{\dot u})ur
   +Y_v v+Y_{vv}|v|v &= 0, \\
 (I_z+N_{\dot r})\dot r +(Y_{\dot v}-X_{\dot u})uv
   +N_r r+N_{rr}|r|r &= T_r.
\end{align}

\section{Comparación actual}
Con
\begin{equation}
m_{11}=m+X_{\dot u},\qquad m_{22}=m+Y_{\dot v},\qquad
m_{33}=I_z+N_{\dot r},
\end{equation}
las tres ecuaciones del simulador corregido son idénticas al modelo estándar de Fossen. Ya no existe una diferencia en los términos $ur$ de sway ni $uv$ de yaw. El identificador de 9 parámetros usa ahora directamente este modelo estándar.

La masa rígida usada es
\begin{equation}
m=21.92+4(0.344+0.012)=23.344\,\mathrm{kg},
\end{equation}
y el momento rígido de guiñada usado es $I_z=2.923312\,\mathrm{kg\,m^2}$.

\section{Resultado del experimento anterior}
Estos valores fueron obtenidos antes de corregir el plugin y se conservan solo como referencia; \textbf{no deben utilizarse con el simulador corregido}. Es necesario ejecutar una nueva excitación y volver a identificar. Con \texttt{wamvsim\_20260908\_020709.mat}, la regresión anterior produjo:
\begin{center}
\begin{tabular}{@{}lr@{}}
\toprule
Parámetro & Valor \\
\midrule
$X_{\dot u}$ & $2.015371\,\mathrm{kg}$ \\
$Y_{\dot v}$ & $7.107892\,\mathrm{kg}$ \\
$N_{\dot r}$ & $2.296427\,\mathrm{kg\,m^2}$ \\
$m_{11}=m+X_{\dot u}$ & $25.359371\,\mathrm{kg}$ \\
$m_{22}=m+Y_{\dot v}$ & $30.451892\,\mathrm{kg}$ \\
$m_{33}=I_z+N_{\dot r}$ & $5.219739\,\mathrm{kg\,m^2}$ \\
$X_u$, $X_{uu}$ & $15.673656$, $14.564905$ \\
$Y_v$, $Y_{vv}$ & $35.470464$, $1.150910$ \\
$N_r$, $N_{rr}$ & $9.367915$, $0.665270$ \\
\bottomrule
\end{tabular}
\end{center}

La próxima identificación, obtenida con datos posteriores a la recompilación, será directamente compatible con Fossen y con el plugin.

\end{document}
